\documentclass[11pt, a4paper]{article}
\usepackage{graphicx}
\usepackage[T2A]{fontenc}
\usepackage[utf8]{inputenc}
\usepackage[russian]{babel}
\usepackage{amsmath, amssymb}
\usepackage[usenames]{color}
\usepackage{colortbl}
\usepackage{fancyhdr}
\usepackage{enumitem}
\usepackage{titlesec}
\usepackage{multicol}
\usepackage{tabularx}
\usepackage{caption}
\usepackage{tikz}
\usetikzlibrary{positioning, arrows.meta}
\usepackage[backend=biber, style=numeric, sorting=none]{biblatex}
\DefineBibliographyExtras{russian}{\protected\def\bibrangedash{\textendash\penalty\hyphenpenalty}}
\AtBeginBibliography{\sloppy}
\renewcommand*{\bibfont}{\small}
\addbibresource{source.bib}
\usepackage{csquotes}
\usepackage[hidelinks]{hyperref}
\usepackage{xurl}
\usepackage{placeins}
\usepackage{longtable}
\usepackage{array}
\usepackage{pdflscape}

\usepackage{geometry}
\geometry{
  a4paper,
  top=30mm,
  right=5mm,
  bottom=20mm,
  left=5mm,
  headheight=61pt
}
\setlength{\columnsep}{8mm}
\raggedcolumns  % не растягивать колонки по высоте
\emergencystretch=2em
\renewcommand{\topfraction}{0.95}
\renewcommand{\bottomfraction}{0.6}
\renewcommand{\textfraction}{0.05}
\renewcommand{\floatpagefraction}{0.85}
\setcounter{topnumber}{3}\setcounter{totalnumber}{4}

\newcommand{\mydate}{Октябрь, 2026}
\newcommand{\myauthor}{В. Алексеев, С. Зинина, И. Коновалов, В. Горбатенко, А. Соколова}

\fancyhf{}
\fancyhead[L]{\large \textit{Научная студия 2026. Вычислительная онкология. Общий анализ данных.}}
\fancyfoot[C]{\thepage}
\pagestyle{fancy}

\fancypagestyle{firstpage}{
  \fancyhf{}
  \fancyhead[L]{\hspace{-15pt}\includegraphics[width=0.3\linewidth]{fig/cu_logo.png}}
  \fancyhead[R]{\mydate\\\textit{\myauthor}\\}
  \fancyfoot[C]{\thepage}
}

\newcommand{\code}[1]{\texttt{\small #1}}

\begin{document}
\thispagestyle{firstpage}

\begin{multicols}{2}
\section*{Материалы и методы}

\subsection*{Исход и момент применения модели}

В наборе два возможных исхода: локальный рецидив первичной опухоли и жизненный статус на момент последнего наблюдения. Жизненный статус известен для всех 298 пациентов, но без времени наблюдения смерть в первый год и через несколько лет неразличимы, а общая выживаемость зависит и от причин, не связанных с облучённым очагом. Поэтому исходом выбран локальный рецидив, заданный авторами набора \cite{Elhalawani2017}: 1~--- рецидив в той же или иной подобласти ротоглотки, 0~--- его отсутствие. Этот исход используется в соревновании на Kaggle \cite{KaggleOPC}, которое принимает решения в режиме поздней отправки (Late Submission), что позволяет независимо проверить предсказания для тестовой части. По условиям соревнования модель строится по первичной опухоли, без признаков лимфоузлов.\vspace{6pt}\\
Исход означает рецидив за всё время наблюдения. Времени до рецидива и длительности наблюдения в таблице нет, поэтому горизонт прогноза задать нельзя. По описанию набора всех пациентов наблюдали более двух лет \cite{Elhalawani2017}, так что отсутствие рецидива означает его отсутствие как минимум за этот срок.\vspace{6pt}\\
Исход известен для 140 пациентов обучающей части: 12 с рецидивом и 128 без него. 158 пациентов тестовой части без значения исхода в аналитическую таблицу не вошли и оставлены для проверки на Kaggle. Малое число событий требует жёсткого ограничения числа признаков при моделировании \cite{Riley2020}.\vspace{6pt}\\
Модель применяется после стадирования и утверждения плана лечения, до его начала: тогда оценка риска позволяет обсуждать эскалацию дозы или более частое наблюдение. К этому моменту известны диагностическая компьютерная томография (КТ), демографические данные, курение, TNM и стадия, степень дифференцировки, назначенные доза и фракционирование и решение о химиотерапии. Поэтому признаками не могут быть фактическая длительность курса лучевой терапии (ЛТ), известная лишь после его окончания, и жизненный статус: это была бы утечка информации из будущего \cite{Kaufman2012}. Доза и число фракций в описании набора приведены как полученные пациентом \cite{Elhalawani2017}, а они известны только после курса. Поэтому принято допущение, что полученные значения совпадают с назначенными: все сочетания дозы и числа фракций соответствуют стандартным схемам (60--72~Гр за 30--40 фракций), а прерванный курс дал бы промежуточные значения. То же допущение принято для химиотерапии.

\subsection*{Первичный анализ данных}

Использована итоговая таблица этапа извлечения признаков: 298 пациентов, идентификатор, часть набора, 18 клинических переменных, число поражённых лимфоузлов и по 107 радиомических признаков макроскопического объёма первичной опухоли (gross tumor volume, GTVp) и лимфоузлов (GTVn); исходный файл не изменялся. Выведены все значения 14 категориальных переменных и наименьшее и наибольшее значения количественных, для клинических~--- в сравнении с допустимым диапазоном. Разных написаний одной категории нет. Пропуски записаны пустыми ячейками (<<не поддаётся оценке>> \cite{Elhalawani2017}), в категории N и стадии есть служебный код X (<<не может быть оценена>>).\vspace{6pt}\\
Значение заменялось верным, только если оно однозначно восстанавливается по другим полям записи, а служебные коды и невосстановимые ошибки заменялись пропуском. Код X у одного пациента заменён пропуском в обеих переменных. У одного пациента при T4N1 указана стадия III, хотя по 7-му изданию классификации Американского объединённого комитета по раку (AJCC) T4 без отдалённых метастазов означает стадию IV \cite{Edge2010}, и стадия исправлена. У 7 курящих и бывших курильщиков указано 0 пачко-лет, что противоречит статусу, и значения заменены пропуском. Всего в клинических переменных изменено 10 значений: 2 в стадии, 1 в категории N и 7 в пачко-годах. Доза за фракцию (1,8--2,4~Гр) и число суток на фракцию (0,98--1,70, меньше единицы при облучении дважды в день) находятся в допустимом диапазоне.\vspace{6pt}\\
Максимальная плотность в маске опухоли выше 3071 единицы Хаунсфилда (HU), верхней границы стандартной шкалы КТ, найдена у 2 пациентов, а у пациента 146 обучающей части достигает 12\,356~HU. У тканей таких значений не бывает, их дают металлические предметы, например зубные пломбы. Металл искажает распределение плотности и текстуру, а значения для ткани не восстановить, поэтому у этих пациентов все 93 признака плотности и текстуры заменены пропуском, признаки формы сохранены. Минимальная плотность ниже $-1024$~HU (ниже воздуха) ещё у 8 пациентов вызвана выбросами интерполяции при передискретизации в отдельных вокселях, и у них пропуском заменены только минимум и размах. Пациенты 219, 257 и 289 с некорректными масками опухоли относятся к тестовой части и в аналитическую таблицу не входят. Для Kaggle нужно предсказание для каждого тестового пациента, поэтому они сохранены, а влияние ошибок разметки будет учтено при интерпретации.\vspace{6pt}\\
Пропуски есть в степени дифференцировки (43), пачко-годах (21: 14 исходных и 7 после исправления), расе, категории N и стадии (по 1), во всех признаках лимфоузлов у 45 пациентов без поражённых узлов, а после исправления~--- в минимуме и размахе плотности (по 10) и остальных признаках плотности и текстуры (по 2). В аналитической когорте их доля составляет 11,4~\% в степени дифференцировки, 7,1~\% в пачко-годах, 5,0~\% в минимуме и размахе и 0,7~\% (один пациент) в остальных признаках плотности и текстуры. Большая доля пропусков считалась доводом за исключение переменной, а не автоматическим критерием. Здесь доли невелики, а степень дифференцировки и курение клинически значимы, поэтому из-за пропусков ничего не исключено.\vspace{6pt}\\
Исключены идентификатор (сохранён как ключ строки), исход, отметка части набора, длительность курса ЛТ, жизненный статус, все 108 признаков лимфоузлов (107 радиомических и число узлов) и дубликат: полная энергия (TotalEnergy) совпадает с энергией (Energy), так как при вокселе 1~мм$^3$ они отличаются лишь множителем \cite{PyRadiomicsDocs}. Констант нет. Добавлена вычисленная доза за фракцию, поскольку при равной суммарной дозе эффект облучения зависит от фракционирования.\vspace{6pt}\\
Избыточность 111 количественных признаков (5 клинических и 106 признаков опухоли) сокращалась последовательным отбором: признак оставлялся, если модуль коэффициента корреляции с каждым ранее оставленным не превышал 0,7, иначе исключался и засчитывался наиболее связанному с ним оставленному. Порядок предпочтения не зависит от исхода: клинические признаки как доступные в рутинной практике, затем объём опухоли (установленный прогностический фактор) и сферичность, далее остальные признаки в порядке PyRadiomics. Использовалась ранговая корреляция Спирмена, так как у 69~\% радиомических признаков модуль коэффициента асимметрии больше 1. Порог 0,7 (почти половина общей ранговой дисперсии) выбран жёстким из-за малого числа событий. Корреляции считались по всем 298 пациентам без использования исхода. Оставлено 20 признаков (4 клинических и 16 радиомических), исключён 91, наибольший модуль корреляции среди оставленных равен 0,674. Больше всего признаков заменяют энтропия (27), минимальная плотность (24) и объём опухоли (23).\vspace{6pt}\\
Аналитическая таблица содержит 140 пациентов, код пациента, исход и 31 признак: 11 категориальных и 20 количественных. Пропуски (16 в степени дифференцировки, 10 в пачко-годах, 7 в минимальной плотности и по одному в остальных 11 признаках плотности и текстуры) на этом этапе не заполнялись, а связь признаков с исходом не изучалась. Заполнение и отбор по исходу будут выполняться внутри перекрёстной проверки, чтобы оценка качества не была оптимистично смещена.

\end{multicols}

\FloatBarrier
\begin{multicols}{2}
\subsection*{Ссылки}
\printbibliography[heading=none]
\end{multicols}

\newpage
\begin{landscape}
\subsection*{Словарь переменных и решений}

\noindent Приведены все клинические и служебные переменные исходной таблицы, вычисленная доза за фракцию, исход и оставленные радиомические признаки первичной опухоли (16 из 107) с числом исключённых признаков, которые каждый из них заменяет. Признаки лимфоузлов исключены целиком: число узлов приведено как переменная исходной таблицы с решением об исключении, а 107 радиомических признаков лимфоузлов в словаре не приводятся. Пропуски указаны после исправления значений во всей таблице (298 пациентов) и в аналитической когорте (140 пациентов). Момент доступности дан относительно применения модели после утверждения плана лечения.

\footnotesize
\setlength{\tabcolsep}{3.5pt}
\renewcommand{\arraystretch}{1.15}
\begin{longtable}{>{\raggedright\arraybackslash}p{5.0cm} >{\raggedright\arraybackslash}p{4.4cm} >{\raggedright\arraybackslash}p{3.6cm} >{\raggedright\arraybackslash}p{3.1cm} >{\raggedright\arraybackslash}p{2.2cm} >{\raggedright\arraybackslash}p{4.9cm}}
\hline\hline
Переменная & Смысл & Тип, допустимые значения & Момент доступности & Пропуски (298 / 140) & Решение и комментарий \\ \hline
\endfirsthead
\hline\hline
Переменная & Смысл & Тип, допустимые значения & Момент доступности & Пропуски (298 / 140) & Решение и комментарий \\ \hline
\endhead
\hline
\multicolumn{6}{r}{\textit{Продолжение на следующей странице}}\\
\endfoot
\hline\hline
\caption{Словарь переменных и решений первичного анализа.}
\endlastfoot
\texttt{Patient\_\allowbreak{}ID} & код пациента, присвоенный при анонимизации, совпадает с именем папки снимков & целый, 1–315, уникален & служебная, всегда & 0 / 0 & оставить как ключ строки, не признак: идентификатор не несёт клинической информации \\
\texttt{Fold} & часть набора соревнования & категориальный: Training, Test & служебная, задана организаторами & 0 / 0 & исключить: служебная переменная \\
\texttt{Local\_\allowbreak{}tumor\_\allowbreak{}recurrence} & локальный рецидив первичной опухоли за период наблюдения & бинарный: 0, 1 & после наблюдения не менее 2 лет & 158 (53,0~\%) / 0 & исключить из признаков: источник целевой переменной \\
\texttt{Gender} & пол & категориальный: Male, Female & до начала лечения, при постановке диагноза & 0 / 0 & оставить \\
\texttt{Age\_\allowbreak{}at\_\allowbreak{}diagnosis} & возраст на момент диагноза, лет & целый, 30–89 & до начала лечения, при постановке диагноза & 0 / 0 & оставить, заменяет 0 исключённых признаков \\
\texttt{Race} & раса & категориальный: White, Black, Hispanic, Asian, American\_Indian/\allowbreak{}Alaska\_Native & до начала лечения, при постановке диагноза & 1 (0,3~\%) / 0 & оставить, пустая ячейка~--- пропуск, не заполнять \\
\texttt{Tumor\_\allowbreak{}side} & латеральность опухоли & категориальный: L, R, Bilateral & до начала лечения, после стадирования & 0 / 0 & оставить \\
\texttt{Tumor\_\allowbreak{}subsite} & подобласть ротоглотки & категориальный: BOT, Tonsil, GPS, Soft\_palate, Pharyngeal\_wall, Other & до начала лечения, после стадирования & 0 / 0 & оставить \\
\texttt{T\_\allowbreak{}category} & категория T (AJCC 7) & порядковый: 1–4 & до начала лечения, после стадирования & 0 / 0 & оставить \\
\texttt{N\_\allowbreak{}category} & категория N (AJCC 7) & порядковый: 0, 1, 2a, 2b, 2c, 3 & до начала лечения, после стадирования & 1 (0,3~\%) / 0 & оставить, служебный код X (1 значение) заменён пропуском \\
\texttt{AJCC\_\allowbreak{}Stage} & стадия по AJCC 7, определяется T и N при M0 & порядковый: I–IV & до начала лечения, после стадирования & 1 (0,3~\%) / 0 & оставить, X (1) заменён пропуском, III при T4N1 (1) исправлено на IV \\
\texttt{Pathological\_\allowbreak{}grade} & степень дифференцировки опухоли & порядковый: I, I-II, II, II-III, III, IV & до начала лечения, по биопсии & 43 (14,4~\%) / 16 (11,4~\%) & оставить: клинически значимый фактор, пропуски на этом этапе не заполнять \\
\texttt{Smoking\_\allowbreak{}status\_\allowbreak{}at\_\allowbreak{}diagnosis} & статус курения на момент диагноза & категориальный: Never, Former, Current & до начала лечения, при постановке диагноза & 0 / 0 & оставить \\
\texttt{Smoking\_\allowbreak{}Pack-Years} & кумулятивное курение, пачко-лет & непрерывный, $\geq 0$, 0 только при Never & до начала лечения, при постановке диагноза & 21 (7,0~\%) / 10 (7,1~\%) & оставить, 7 нулей при Current/Former заменены пропуском, заменяет 0 исключённых признаков \\
\texttt{Radiation\_\allowbreak{}treatment\_\allowbreak{}course\_\allowbreak{}duration} & фактическая длительность курса ЛТ, дни & целый, 32–56 & после окончания ЛТ & 0 / 0 & исключить: неизвестна в момент применения модели \\
\texttt{Total\_\allowbreak{}prescribed\_\allowbreak{}Radiation\_\allowbreak{}treatment\_\allowbreak{}dose} & суммарная доза, Гр (в описании набора~--- полученная) & целый: 60, 63, 66, 70, 72 & до начала ЛТ как назначенное значение (допущение: получено как назначено) & 0 / 0 & оставить: все значения~--- стандартные назначенные дозы, поэтому полученная доза принята равной назначенной, заменяет 0 исключённых признаков \\
\texttt{\#\_\allowbreak{}Radiation\_\allowbreak{}treatment\_\allowbreak{}fractions} & число фракций (в описании набора~--- полученное) & целый, 30–40 & до начала ЛТ как назначенное значение (допущение: получено как назначено) & 0 / 0 & исключить: избыточна, $|\rho| > 0{,}7$ с Dose\_per\_fraction \\
\texttt{Induction\_\allowbreak{}Chemotherapy} & индукционная химиотерапия до ЛТ & бинарный: Y, N & до начала ЛТ как назначенное значение (допущение: получено как назначено) & 0 / 0 & оставить: принято, что проведена как запланирована \\
\texttt{Concurrent\_\allowbreak{}chemotherapy} & химиотерапия одновременно с ЛТ & бинарный: Y, N & до начала ЛТ как назначенное значение (допущение: получено как назначено) & 0 / 0 & оставить: принято, что проведена как запланирована \\
\texttt{KM\_\allowbreak{}Overall\_\allowbreak{}survival\_\allowbreak{}censor} & жизненный статус на последнем наблюдении & бинарный: 1~--- жив, 0~--- умер & после наблюдения & 0 / 0 & исключить: альтернативный исход, известен в обеих частях набора, но без времени наблюдения, неизвестен в момент применения модели \\
\texttt{GTVn\_\allowbreak{}n\_\allowbreak{}nodes} & число поражённых лимфоузлов, выделенных на маске GTVn & целый, 0–16 & до начала лечения, по диагностической КТ & 0 / 0 & исключить вместе со всеми признаками GTVn: по условиям соревнования модель строится по первичной опухоли \\
\texttt{Dose\_\allowbreak{}per\_\allowbreak{}fraction} & вычисленная: доза за фракцию, Гр (доза / число фракций) & непрерывный, 1,8–2,4 & до начала ЛТ как назначенное значение (допущение: получено как назначено) & 0 / 0 & оставить, заменяет 1 исключённый признак \\
\texttt{Target\_\allowbreak{}local\_\allowbreak{}recurrence} & целевая: 1~--- локальный рецидив, 0~--- нет & бинарный: 0 (128), 1 (12) & после наблюдения не менее 2 лет & 158 (53,0~\%) / 0 & целевая переменная \\
\texttt{GTVp\_\allowbreak{}shape\_\allowbreak{}Mesh\allowbreak{}Volume} & объём по треугольной сетке (GTVp) & непрерывный, мм$^3$, $>0$ & до начала лечения, по диагностической КТ & 0 / 0 & оставить, заменяет 23 исключённых признака \\
\texttt{GTVp\_\allowbreak{}shape\_\allowbreak{}Sphericity} & сферичность (GTVp) & непрерывный, (0, 1] & до начала лечения, по диагностической КТ & 0 / 0 & оставить, заменяет 0 исключённых признаков \\
\texttt{GTVp\_\allowbreak{}shape\_\allowbreak{}Elongation} & вытянутость: отношение второй главной оси к первой (GTVp) & непрерывный, (0, 1] & до начала лечения, по диагностической КТ & 0 / 0 & оставить, заменяет 0 исключённых признаков \\
\texttt{GTVp\_\allowbreak{}shape\_\allowbreak{}Flatness} & уплощённость: отношение третьей главной оси к первой (GTVp) & непрерывный, (0, 1] & до начала лечения, по диагностической КТ & 0 / 0 & оставить, заменяет 0 исключённых признаков \\
\texttt{GTVp\_\allowbreak{}firstorder\_\allowbreak{}10Percentile} & 10-й перцентиль плотности (GTVp) & непрерывный, HU & до начала лечения, по диагностической КТ & 2 (0,7~\%) / 1 (0,7~\%) & оставить, заменяет 1 исключённый признак, у 2 пациентов с металлом в маске значение заменено пропуском \\
\texttt{GTVp\_\allowbreak{}firstorder\_\allowbreak{}90Percentile} & 90-й перцентиль плотности (GTVp) & непрерывный, HU & до начала лечения, по диагностической КТ & 2 (0,7~\%) / 1 (0,7~\%) & оставить, заменяет 2 исключённых признака, у 2 пациентов с металлом в маске значение заменено пропуском \\
\texttt{GTVp\_\allowbreak{}firstorder\_\allowbreak{}Entropy} & энтропия гистограммы плотностей (GTVp) & непрерывный, бит, $\geq 0$ & до начала лечения, по диагностической КТ & 2 (0,7~\%) / 1 (0,7~\%) & оставить, заменяет 27 исключённых признаков, у 2 пациентов с металлом в маске значение заменено пропуском \\
\texttt{GTVp\_\allowbreak{}firstorder\_\allowbreak{}Kurtosis} & эксцесс распределения плотностей (GTVp) & непрерывный, $\geq 1$ & до начала лечения, по диагностической КТ & 2 (0,7~\%) / 1 (0,7~\%) & оставить, заменяет 3 исключённых признака, у 2 пациентов с металлом в маске значение заменено пропуском \\
\texttt{GTVp\_\allowbreak{}firstorder\_\allowbreak{}Maximum} & максимальная плотность (GTVp) & непрерывный, HU, от $-1024$ до 3071 & до начала лечения, по диагностической КТ & 2 (0,7~\%) / 1 (0,7~\%) & оставить, заменяет 0 исключённых признаков, у 2 пациентов с металлом в маске значение заменено пропуском \\
\texttt{GTVp\_\allowbreak{}firstorder\_\allowbreak{}Minimum} & минимальная плотность (GTVp) & непрерывный, HU, от $-1024$ до 3071 & до начала лечения, по диагностической КТ & 10 (3,4~\%) / 7 (5,0~\%) & оставить, заменяет 24 исключённых признака, у 2 пациентов с металлом в маске значение заменено пропуском, ещё 8 значений ниже $-1024$~HU заменены пропуском \\
\texttt{GTVp\_\allowbreak{}glcm\_\allowbreak{}Correlation} & GLCM: корреляция уровней соседних вокселей (GTVp) & непрерывный, $[-1, 1]$ & до начала лечения, по диагностической КТ & 2 (0,7~\%) / 1 (0,7~\%) & оставить, заменяет 4 исключённых признака, у 2 пациентов с металлом в маске значение заменено пропуском \\
\texttt{GTVp\_\allowbreak{}glcm\_\allowbreak{}Inverse\allowbreak{}Variance} & GLCM: обратная дисперсия (GTVp) & непрерывный, $\geq 0$ & до начала лечения, по диагностической КТ & 2 (0,7~\%) / 1 (0,7~\%) & оставить, заменяет 0 исключённых признаков, у 2 пациентов с металлом в маске значение заменено пропуском \\
\texttt{GTVp\_\allowbreak{}gldm\_\allowbreak{}Dependence\allowbreak{}Non\allowbreak{}Uniformity\allowbreak{}Normalized} & GLDM: нормированная неоднородность зависимостей (GTVp) & непрерывный, (0, 1] & до начала лечения, по диагностической КТ & 2 (0,7~\%) / 1 (0,7~\%) & оставить, заменяет 5 исключённых признаков, у 2 пациентов с металлом в маске значение заменено пропуском \\
\texttt{GTVp\_\allowbreak{}glszm\_\allowbreak{}Size\allowbreak{}Zone\allowbreak{}Non\allowbreak{}Uniformity\allowbreak{}Normalized} & GLSZM: нормированная неоднородность размеров зон (GTVp) & непрерывный, (0, 1] & до начала лечения, по диагностической КТ & 2 (0,7~\%) / 1 (0,7~\%) & оставить, заменяет 1 исключённый признак, у 2 пациентов с металлом в маске значение заменено пропуском \\
\texttt{GTVp\_\allowbreak{}ngtdm\_\allowbreak{}Busyness} & NGTDM: «занятость»~--- частота смены тонов (GTVp) & непрерывный, $\geq 0$ & до начала лечения, по диагностической КТ & 2 (0,7~\%) / 1 (0,7~\%) & оставить, заменяет 1 исключённый признак, у 2 пациентов с металлом в маске значение заменено пропуском \\
\texttt{GTVp\_\allowbreak{}ngtdm\_\allowbreak{}Contrast} & NGTDM: контраст соседних тонов (GTVp) & непрерывный, $\geq 0$ & до начала лечения, по диагностической КТ & 2 (0,7~\%) / 1 (0,7~\%) & оставить, заменяет 0 исключённых признаков, у 2 пациентов с металлом в маске значение заменено пропуском \\

\end{longtable}
\end{landscape}

\end{document}
